\section*{Sets and Indices}

\begin{itemize}
    \item $R$ : set of clerk roles. In this instance,
    \[
    R=\{\mathrm{FT},\mathrm{PT}\},
    \]
    where $\mathrm{FT}$ denotes full-time clerks and $\mathrm{PT}$ denotes part-time clerks.
\end{itemize}

\section*{Parameters}

\begin{itemize}
    \item $G$ : monthly sales goal in pairs (code: \texttt{sales\_goal}).
    \item $n_r$ : number of clerks of role $r$ employed by the store. In the implemented code, if no employee-count column is present in the CSV, the defaults are
    \[
    n_{\mathrm{FT}}=5,\qquad n_{\mathrm{PT}}=4
    \]
    (code: \texttt{num\_ft}, \texttt{num\_pt}).
    \item $H_r$ : regular monthly working hours per clerk of role $r$ from the CSV column \texttt{Monthly\_Working\_Hours}.
    \item $s_r$ : sales volume in pairs per hour for role $r$ from the CSV column \texttt{Sales\_Volume\_Pairs\_Per\_Hour}.
    \item $w_r$ : regular wage per hour for role $r$ from the CSV column \texttt{Wage\_Yuan\_Per\_Hour}. This parameter is read from data but not used in the implemented optimization model.
    \item $o_r$ : overtime pay per hour for role $r$ from the CSV column \texttt{Overtime\_Pay\_Yuan\_Per\_Hour}. This parameter is read from data but not used in the implemented optimization model.
    \item $p$ : profit per pair sold from \texttt{profit\_per\_pair}. This parameter is not used in the implemented optimization model.
    \item $N$ : total sales achievable under full regular employment, computed as
    \[
    N=\sum_{r\in R} n_r H_r s_r
    \]
    (code: \texttt{total\_normal\_sales}).
\end{itemize}

For this instance, using the provided table values,
\[
H_{\mathrm{FT}}=160,\quad s_{\mathrm{FT}}=5,\qquad
H_{\mathrm{PT}}=80,\quad s_{\mathrm{PT}}=2,
\]
so
\[
N=n_{\mathrm{FT}} \cdot 160 \cdot 5 + n_{\mathrm{PT}} \cdot 80 \cdot 2.
\]

\section*{Decision Variables}

\begin{itemize}
    \item $y_r \ge 0$ : overtime hours assigned to role $r$ in aggregate over the month (code: \texttt{y\_ft}, \texttt{y\_pt} for $r=\mathrm{FT},\mathrm{PT}$).
\end{itemize}

Explicitly,
\[
y_{\mathrm{FT}} \ge 0,\qquad y_{\mathrm{PT}} \ge 0.
\]

\section*{Objective}

Minimize total overtime hours:
\[
\min \; y_{\mathrm{FT}} + y_{\mathrm{PT}}.
\]

\section*{Constraints}

Monthly sales goal must be met after accounting for regular-hours sales and overtime sales:
\[
s_{\mathrm{FT}}\, y_{\mathrm{FT}} + s_{\mathrm{PT}}\, y_{\mathrm{PT}} \;\ge\; G - N.
\]

Nonnegativity:
\[
y_{\mathrm{FT}} \ge 0,\qquad y_{\mathrm{PT}} \ge 0.
\]

Using the instance data and the code defaults $n_{\mathrm{FT}}=5$, $n_{\mathrm{PT}}=4$,
\[
N=5\cdot160\cdot5+4\cdot80\cdot2=4640,
\]
and with $G=5500$ the implemented constraint becomes
\[
5y_{\mathrm{FT}}+2y_{\mathrm{PT}} \ge 860.
\]

\section*{Notes on Implementation}

\begin{itemize}
    \item The model solved in the code is exactly a continuous linear program with only two decision variables: aggregate overtime hours for full-time and part-time clerks.
    \item The code assumes ``full employment'' by always counting all available full-time and part-time clerks in the baseline sales computation $N$. No decision variable controls hiring, assignment, or whether clerks are employed.
    \item Although the business description mentions profit, wages, overtime pay, and multiple goals, the implemented model ignores all of these except the monthly sales goal and the role-specific sales rates. In particular, \texttt{profit\_per\_pair}, regular wages, and overtime pay do not appear in either the objective or the constraints.
    \item Overtime is modeled in aggregate by role, not per individual clerk, and there are no upper bounds on overtime hours.
    \item If the CSV contains a column whose name suggests employee counts (matching keywords such as \texttt{num}, \texttt{clerks}, \texttt{employees}, \texttt{count}, etc.), the code uses those values for $n_{\mathrm{FT}}$ and $n_{\mathrm{PT}}$; otherwise it uses the defaults $5$ and $4$.
    \item The LP is solved with a standard LP/MIP solver interface (\texttt{GUROBI\_CMD} through PuLP compatibility), not by enumeration or dynamic programming.
\end{itemize}
